\section{UV divergences at high orders}
\label{sec:top}
From the one-loop diagrams in Fig.~\ref{fig:oneloop-g}, we can generate all high order loop diagrams by adding gluons, quark-antiquark pairs, or ghost-antighost pairs to them. Because of the isolation of $z$-component in the definition of quasi-parton operators, both 3-dimensional (3-D) and 4-D integration of loop momentum $\bar{l}$ and $l$ could lead to UV divergence.  In Ref.~\cite{Ishikawa:2017faj}, we introduced the change of divergence index $\Delta\omega_3$ and $\Delta\omega_4$ for the 3-D and 4-D integration of loop momenta of higher order diagrams, respectively, and showed that it is sufficient, although it is not necessary, that quasi-parton operators are renormalizable if $\Delta\omega_3\leq 0$ and $\Delta\omega_4\leq 0$ are satisfied for all corresponding higher order diagrams.  Based on the power counting rules derived in Ref.~\cite{Ishikawa:2017faj}, we find that the only case that may increase superficial UV divergence of quasi-gluon operators at high orders is when we add a gluon with both ends of it attached to the gauge link, where the 3-D integration gets $\Delta\omega_3=1$. By applying the decomposition in Eq.~\eqref{eq:dec} to the added gluon's momentum, it is straightforward to show that dimensional regularized UV divergences at any loop level are localized in spacetime, in the same way as quasi-quark operators shown in~\cite{Ishikawa:2017faj}.  As a result, we find that $\Delta\omega_3$ is effectively irrelevant for studying UV divergences. Because $\Delta\omega_4\leq0$ for all cases, there are only finite topologies of high order diagrams in Fig.~\ref{fig:div-topo-g} that have UV divergences.

%%%%%%%%%%%%%%%%%%%%%%%%%%%%%%%
\begin{figure}[htb]
	\begin{center}
		\includegraphics[width=0.75\textwidth]{images/div-topo-g.pdf}
		\caption{\label{fig:div-topo-g}
			Four topologies of diagrams which could give UV divergences to the quasi-gluon operators.}
	\end{center}
\end{figure}
%%%%%%%%%%%%%%%%%%%%%%%%%%%%%%

The blobs with topologies (a) and (c) in Fig.~\ref{fig:div-topo-g} denote one-particle-irreducible diagrams, and they both have linear superficial UV divergences. Because of the potential linear UV divergences, diagrams with one more gluon attached to the blobs can generate logarithmic UV divergences.  Another possibility to produce logarithmic divergences is when a gluon is attached to the gauge link outside of, but very close to, the blobs, as shown in Fig.~\ref{fig:div-topo-g} (b) and (d), with the attachment denoted by a triangle. The blobs of topologies (b) and (d) include both kinds of logarithmic divergent diagrams mentioned here.

The topologies (a) and (b) in Fig.~\ref{fig:div-topo-g} are the same as that for quasi-quark operators, and their divergences can be renormalized similarly.  Linear divergences from the diagrams of the topology (a) can be removed by an overall factor as the mass renormalization of a test particle moving along the gauge link \cite{Polyakov:1980ca}, and its logarithmic divergences caused by end points of the gauge link can be removed by multiplying $Z_{wg}^{-1/2}$ - the ``wave function" renormalization of the test particle \cite{Dotsenko:1979wb}. The diagrams of the topology (b) has only logarithmic UV divergence, which can be taken care of by QCD  renormalization \cite{Dotsenko:1979wb}.

The UV divergences from diagrams of topologies (c) and (d) in Fig.~\ref{fig:div-topo-g} are different from that of quasi-quark operators, and are studied in next two sections, respectively.
